\section{Población}

La población de estudio se define como un conjunto cerrado, finito y accesible de casos (personas, animales, expedientes, objetos, etc.), los cuales cumplen con un criterio de inclusión y que sirven como base para la investigación \cite{poblacion}. En otras palabras, es un conjunto conformado por todos los elementos a los que se pretende generalizar los resultados del estudio.

En este trabajo, la población que se estudia está formada por todos los autómatas celulares finitos que se encuentran en un espacio bidimensional con una distribución de cuadrícula regular. Se consideran condiciones de frontera un conjunto de estados finito y cualquier función que pueda mapear de $\mathbb{R}^{|\mathcal{N}|} \to \mathbb{R}$.

